\chapter{Perspectives et Améliorations}

\section{Introduction}

LudoHeritage est une plateforme fonctionnelle qui répond aux objectifs fixés dans le cadre de
ce projet de fin d'année. Cependant, toute application a une marge de progression. Ce chapitre
présente les améliorations envisagées pour faire évoluer la plateforme vers une solution encore
plus complète et robuste.

% ─────────────────────────────────────────────────────────────────────────────
\section{Améliorations techniques}
% ─────────────────────────────────────────────────────────────────────────────

\subsection{Authentification sécurisée par tokens JWT}

L'implémentation actuelle utilise le stockage de la session utilisateur dans \texttt{localStorage}.
Pour renforcer la sécurité, une prochaine version pourrait intégrer des \textbf{tokens JWT}
(JSON Web Tokens) avec un mécanisme de rafraîchissement automatique. Cela permettrait une
authentification \textit{stateless} plus robuste et une meilleure gestion des sessions.

\subsection{API de jeux enrichie}

Le catalogue actuel est statique et géré entièrement par le code Java. Une amélioration
naturelle serait de créer une \textbf{interface d'administration} permettant à un gestionnaire
de contenu d'ajouter, modifier ou supprimer des jeux directement depuis une interface web,
sans intervention technique.

\subsection{Déploiement en production}

La plateforme tourne actuellement en environnement local. Un déploiement sur le cloud
(par exemple avec \textbf{Docker} + \textbf{Railway} ou \textbf{Render}) permettrait de la
rendre accessible au public. Le modèle Ollama pourrait être remplacé par une API cloud
(Anthropic, Mistral AI) pour les environnements sans GPU.

% ─────────────────────────────────────────────────────────────────────────────
\section{Améliorations fonctionnelles}
% ─────────────────────────────────────────────────────────────────────────────

\subsection{Jeux jouables additionnels}

Seul l'Awale est actuellement jouable dans le navigateur. Des jeux comme le \textbf{Moulin}
(Nine Men's Morris) ou le \textbf{Mancala} pourraient être implémentés avec des interfaces
visuelles dédiées, enrichissant considérablement l'expérience éducative.

\subsection{Multilingue complet}

L'application supporte déjà trois langues dans les préférences utilisateur. Une traduction
complète de l'interface et du contenu des fiches de jeux en anglais et en arabe élargirait
le public cible et renforcerait la dimension multiculturelle du projet.

\subsection{Système de contributions communautaires}

Les utilisateurs passionnés pourraient proposer de nouvelles variantes de règles, partager
des ressources historiques ou suggérer de nouveaux jeux à intégrer au catalogue. Un système
de modération permettrait de valider ces contributions avant publication.

\subsection{LudoBot amélioré}

LudoBot pourrait être enrichi avec~:
\begin{itemize}
    \item Une \textbf{mémoire de conversation} pour maintenir le contexte sur plusieurs échanges.
    \item Un \textbf{mode multilingue}~: détecter automatiquement la langue de l'utilisateur
          et répondre dans cette langue.
    \item Des \textbf{réponses enrichies} incluant des liens cliquables vers les fiches de jeux
          mentionnés.
\end{itemize}

% ─────────────────────────────────────────────────────────────────────────────
\section{Extension du catalogue}
% ─────────────────────────────────────────────────────────────────────────────

Le catalogue actuel de vingt jeux est volontairement limité pour garantir la qualité et
l'exactitude des informations. À terme, LudoHeritage pourrait s'étendre à~:

\begin{itemize}
    \item \textbf{50 à 100 jeux}, en s'appuyant sur des sources académiques fiables comme
          les publications du Digital Ludeme Project.
    \item Des jeux \textbf{régionaux marocains et maghrébins} peu documentés numériquement,
          contribuant à la préservation d'un patrimoine local.
    \item Des jeux avec des \textbf{vidéos explicatives} intégrées pour les règles complexes.
\end{itemize}

% ─────────────────────────────────────────────────────────────────────────────
\section{Conclusion}
% ─────────────────────────────────────────────────────────────────────────────

LudoHeritage constitue une base solide et évolutive. Les perspectives présentées dans ce
chapitre montrent que le projet peut évoluer bien au-delà de sa version actuelle, vers une
plateforme de référence pour la préservation et la diffusion du patrimoine ludique mondial
en langue française. Ces évolutions futures seront guidées par les retours des utilisateurs
et les opportunités offertes par les nouvelles technologies.
